% =====================================================================
%  capitulos/07_anexo_memoria.tex  --  Anexo F: Memoria de calculo
%  Ejemplo Q4 canonico resuelto paso a paso. Los numeros provienen de
%  figuras/anexo_calculo_data.tex (generado por gen_anexo_calculo.py:
%  salida bit-exacta del motor de EduFEM). Las figuras mem_*.png las
%  produce el mismo script con file_io/figure_export (estilo del canvas).
% =====================================================================

% Generado por gen_anexo_calculo.py — NO editar a mano.
% Numeros bit-exactos del motor de EduFEM (ejemplo Q4).

\newcommand{\Bmata}{\begin{bmatrix}
-0{,}1657 & 0{,}0000 & 0{,}2071 & 0{,}0000 & 0{,}0444 & 0{,}0000 & -0{,}0858 & 0{,}0000 \\
0{,}0000 & -0{,}0828 & 0{,}0000 & -0{,}1465 & 0{,}0000 & 0{,}0222 & 0{,}0000 & 0{,}2071 \\
-0{,}0828 & -0{,}1657 & -0{,}1465 & 0{,}2071 & 0{,}0222 & 0{,}0444 & 0{,}2071 & -0{,}0858
\end{bmatrix}}
\newcommand{\Dmat}{\begin{bmatrix}
234375{,}00 & 46875{,}00 & 0{,}00 \\
46875{,}00 & 234375{,}00 & 0{,}00 \\
0{,}00 & 0{,}00 & 93750{,}00
\end{bmatrix}}
\newcommand{\DmatNorm}{\begin{bmatrix}
1{,}00 & 0{,}20 & 0{,}00 \\
0{,}20 & 1{,}00 & 0{,}00 \\
0{,}00 & 0{,}00 & 0{,}40
\end{bmatrix}}
\newcommand{\Eextrap}{\begin{bmatrix}
1{,}8660 & -0{,}5000 & 0{,}1340 & -0{,}5000 \\
-0{,}5000 & 0{,}1340 & -0{,}5000 & 1{,}8660 \\
0{,}1340 & -0{,}5000 & 1{,}8660 & -0{,}5000 \\
-0{,}5000 & 1{,}8660 & -0{,}5000 & 0{,}1340
\end{bmatrix}}
\newcommand{\FfVec}{\begin{bmatrix}
0 \\
0 \\
0 \\
0 \\
0 \\
0 \\
0 \\
-1000 \\
0 \\
0 \\
0 \\
0
\end{bmatrix}}
\newcommand{\Jmata}{\begin{bmatrix}
2{,}1830 & 0{,}3943 \\
1{,}1830 & 2{,}3943
\end{bmatrix}}
\newcommand{\RxTot}{0{,}00}
\newcommand{\RyTot}{1000{,}00}
\newcommand{\Xemat}{\begin{bmatrix}
2{,}0 & 3{,}0 \\
7{,}0 & 4{,}0 \\
7{,}0 & 8{,}0 \\
5{,}0 & 8{,}0
\end{bmatrix}}
\newcommand{\dNnatA}{\begin{bmatrix}
-0{,}3943 & 0{,}3943 & 0{,}1057 & -0{,}1057 \\
-0{,}3943 & -0{,}1057 & 0{,}1057 & 0{,}3943
\end{bmatrix}}
\newcommand{\dNphysA}{\begin{bmatrix}
-0{,}16568 & 0{,}20709 & 0{,}04439 & -0{,}08581 \\
-0{,}08284 & -0{,}14645 & 0{,}02220 & 0{,}20709
\end{bmatrix}}
\newcommand{\despRows}{1 & $0$ & $0$ \\
2 & $-6{,}1088\times10^{-4}$ & $-4{,}1462\times10^{-3}$ \\
3 & $0$ & $0$ \\
4 & $-3{,}8514\times10^{-4}$ & $-4{,}2722\times10^{-3}$ \\
5 & $-1{,}0780\times10^{-3}$ & $-4{,}3740\times10^{-3}$ \\
6 & $0$ & $0$ \\
7 & $-6{,}8936\times10^{-3}$ & $-1{,}3076\times10^{-2}$ \\
8 & $-5{,}3984\times10^{-3}$ & $-5{,}6052\times10^{-3}$ \\
9 & $-5{,}7844\times10^{-3}$ & $-2{,}3686\times10^{-3}$}
\newcommand{\detJa}{4{,}7604}
\newcommand{\detJinline}{4{,}7604,\ 3{,}0283,\ 4{,}4717,\ 2{,}7396}
\newcommand{\detJrow}{4{,}7604 & 3{,}0283 & 4{,}4717 & 2{,}7396}
\newcommand{\elementosRows}{E1 & 1 & 2 & 5 & 4 & $0{,}8$ \\
E2 & 2 & 3 & 6 & 5 & $0{,}8$ \\
E3 & 4 & 5 & 8 & 7 & $0{,}8$ \\
E4 & 5 & 6 & 9 & 8 & $0{,}8$}
\newcommand{\energiaInline}{E1\ $0{,}78$;\ E2\ $1{,}34$;\ E3\ $4{,}35$;\ E4\ $0{,}07$}
\newcommand{\energiaRows}{E1 & $7{,}780\times10^{-1}$ \\
E2 & $1{,}343\times10^{0}$ \\
E3 & $4{,}352\times10^{0}$ \\
E4 & $6{,}514\times10^{-2}$}
\newcommand{\epsA}{\begin{bmatrix}
1{,}924718\times10^{-4} \\
-1{,}837951\times10^{-3} \\
-6{,}824468\times10^{-4}
\end{bmatrix}}
\newcommand{\freeDofList}{2, 3, 6, 7, 8, 9, 12, 13, 14, 15, 16, 17}
\newcommand{\gaussStressRows}{1 & $-41{,}04$ & $-421{,}75$ & $-63{,}98$ & $417{,}76$ \\
2 & $-6{,}76$ & $-611{,}80$ & $87{,}20$ & $626{,}92$ \\
3 & $1{,}82$ & $-167{,}10$ & $-63{,}90$ & $201{,}20$ \\
4 & $66{,}82$ & $-216{,}19$ & $103{,}25$ & $312{,}46$}
\newcommand{\gpA}{\left(-\tfrac{1}{\sqrt{3}},\,-\tfrac{1}{\sqrt{3}}\right)}
\newcommand{\invJa}{\begin{bmatrix}
0{,}50297 & -0{,}08284 \\
-0{,}24851 & 0{,}45858
\end{bmatrix}}
\newcommand{\keMat}{\begin{bmatrix}
63{,}2 & 13{,}5 & -60{,}2 & 1{,}9 & -32{,}5 & -19{,}2 & 29{,}5 & 3{,}9 \\
13{,}5 & 39{,}5 & 20{,}7 & -2{,}6 & -19{,}2 & -13{,}1 & -14{,}9 & -23{,}9 \\
-60{,}2 & 20{,}7 & 122{,}2 & -49{,}3 & 21{,}8 & -13{,}5 & -83{,}8 & 42{,}1 \\
1{,}9 & -2{,}6 & -49{,}3 & 120{,}4 & 5{,}3 & -30{,}5 & 42{,}1 & -87{,}4 \\
-32{,}5 & -19{,}2 & 21{,}8 & 5{,}3 & 117{,}0 & 3{,}3 & -106{,}3 & 10{,}6 \\
-19{,}2 & -13{,}1 & -13{,}5 & -30{,}5 & 3{,}3 & 85{,}8 & 29{,}3 & -42{,}2 \\
29{,}5 & -14{,}9 & -83{,}8 & 42{,}1 & -106{,}3 & 29{,}3 & 160{,}6 & -56{,}5 \\
3{,}9 & -23{,}9 & 42{,}1 & -87{,}4 & 10{,}6 & -42{,}2 & -56{,}5 & 153{,}4
\end{bmatrix}}
\newcommand{\keScale}{10^{3}}
\newcommand{\nDof}{18}
\newcommand{\nFree}{12}
\newcommand{\nRestr}{6}
\newcommand{\nodalAllRows}{1 & $0$ & $0$ & $-57{,}15$ & $-151{,}06$ & $-114{,}37$ & $211{,}90$ \\
2 & $-6{,}109\times10^{-4}$ & $-4{,}146\times10^{-3}$ & $-17{,}82$ & $-118{,}85$ & $73{,}09$ & $131{,}96$ \\
3 & $0$ & $0$ & $34{,}04$ & $-162{,}76$ & $71{,}27$ & $218{,}12$ \\
4 & $-3{,}851\times10^{-4}$ & $-4{,}272\times10^{-3}$ & $-108{,}01$ & $-397{,}47$ & $-86{,}54$ & $305{,}40$ \\
5 & $-1{,}078\times10^{-3}$ & $-4{,}374\times10^{-3}$ & $70{,}30$ & $-54{,}92$ & $50{,}26$ & $187{,}67$ \\
6 & $0$ & $0$ & $-31{,}88$ & $-106{,}15$ & $71{,}23$ & $145{,}90$ \\
7 & $-6{,}894\times10^{-3}$ & $-1{,}308\times10^{-2}$ & $15{,}95$ & $-876{,}18$ & $240{,}48$ & $902{,}22$ \\
8 & $-5{,}398\times10^{-3}$ & $-5{,}605\times10^{-3}$ & $-54{,}23$ & $-11{,}55$ & $-84{,}65$ & $8{,}96$ \\
9 & $-5{,}784\times10^{-3}$ & $-2{,}369\times10^{-3}$ & $104{,}17$ & $1{,}86$ & $-37{,}62$ & $74{,}21$}
\newcommand{\nodalStressShowRows}{4 & $-106{,}37$ & $-395{,}38$ & $-223{,}17$ & $336{,}82$ \\
5 & $143{,}39$ & $-190{,}95$ & $268{,}29$ & $357{,}56$ \\
8 & $-32{,}13$ & $45{,}68$ & $-223{,}04$ & $-38{,}28$ \\
7 & $15{,}95$ & $-876{,}18$ & $240{,}48$ & $902{,}22$}
\newcommand{\nodosRows}{1 & $0{,}0$ & $0{,}0$ \\
2 & $5{,}0$ & $0{,}0$ \\
3 & $8{,}0$ & $0{,}0$ \\
4 & $2{,}0$ & $3{,}0$ \\
5 & $7{,}0$ & $4{,}0$ \\
6 & $11{,}0$ & $4{,}0$ \\
7 & $5{,}0$ & $8{,}0$ \\
8 & $7{,}0$ & $8{,}0$ \\
9 & $9{,}0$ & $8{,}0$}
\newcommand{\radiusA}{200{,}817}
\newcommand{\reaccRows}{1 & $80{,}11$ & $334{,}24$ \\
3 & $-209{,}93$ & $601{,}37$ \\
6 & $129{,}82$ & $64{,}39$}
\newcommand{\restrDofList}{0, 1, 4, 5, 10, 11}
\newcommand{\savgA}{-231{,}395}
\newcommand{\showId}{3}
\newcommand{\showNodes}{4, 5, 8, 7}
\newcommand{\sigA}{\begin{bmatrix}
-41{,}043 \\
-421{,}748 \\
-63{,}979
\end{bmatrix}}
\newcommand{\sigOneA}{-30{,}579}
\newcommand{\sigTwoA}{-432{,}212}
\newcommand{\sigxGauss}{\begin{bmatrix}
-41{,}04 \\
-6{,}76 \\
1{,}82 \\
66{,}82
\end{bmatrix}}
\newcommand{\sigxNodal}{\begin{bmatrix}
-106{,}37 \\
143{,}39 \\
-32{,}13 \\
15{,}95
\end{bmatrix}}
\newcommand{\umaxComp}{uy}
\newcommand{\umaxCompNode}{7}
\newcommand{\umaxCompVal}{1{,}3076\times10^{-2}}
\newcommand{\umaxMag}{1{,}4782\times10^{-2}}
\newcommand{\umaxNode}{7}
\newcommand{\valDfactor}{234375}
\newcommand{\valE}{225000}
\newcommand{\valNu}{0{,}2}
\newcommand{\valP}{1000}
\newcommand{\valT}{0{,}8}
\newcommand{\vmA}{417{,}76}
\newcommand{\vmMax}{902{,}22}
\newcommand{\vmMaxNode}{7}
\newcommand{\vmNodalRows}{1 & $-57{,}15$ & $-151{,}06$ & $-114{,}37$ & $211{,}90$ \\
2 & $-17{,}82$ & $-118{,}85$ & $73{,}09$ & $131{,}96$ \\
3 & $34{,}04$ & $-162{,}76$ & $71{,}27$ & $218{,}12$ \\
4 & $-108{,}01$ & $-397{,}47$ & $-86{,}54$ & $305{,}40$ \\
5 & $70{,}30$ & $-54{,}92$ & $50{,}26$ & $187{,}67$ \\
6 & $-31{,}88$ & $-106{,}15$ & $71{,}23$ & $145{,}90$ \\
7 & $15{,}95$ & $-876{,}18$ & $240{,}48$ & $902{,}22$ \\
8 & $-54{,}23$ & $-11{,}55$ & $-84{,}65$ & $8{,}96$ \\
9 & $104{,}17$ & $1{,}86$ & $-37{,}62$ & $74{,}21$}
\newcommand{\weightsRow}{1 & 1 & 1 & 1}


\chapter{Memoria de cálculo: ejemplo resuelto paso a paso}
\label{anexo:memoria}

Este anexo presenta, en forma condensada, el contenido de la \emph{Memoria de Cálculo} que EduFEM genera para el ejemplo canónico de validación: nueve nodos, cuatro elementos Q4, tensión plana. La memoria completa que produce el software consta de nueve capítulos en dos estilos (\autoref{sec:memoria-io}); aquí se conserva la secuencia de sustitución numérica esencial. Cada matriz, vector y valor que sigue es la salida real del motor numérico, generada con las mismas funciones que emplea EduFEM, con sustitución directa sobre la formulación matricial. Junto a cada expresión se remite a la ecuación numerada del \autoref{cap:marco_teorico} de la que procede, de modo que todo paso numérico pueda cotejarse con su definición simbólica. Las magnitudes se expresan en un sistema de unidades coherente y arbitrario (por ejemplo, kN y cm, con $E$ en kN/cm$^2$): la formulación no depende de la elección mientras sea coherente.

\section{Datos}

La \autoref{fig:mem-modelo} muestra el modelo discretizado; las Tablas \ref{tab:mem-nodos} y \ref{tab:mem-elem} recogen sus nodos y elementos, y la \autoref{tab:mem-param} los parámetros globales del análisis.

\begin{figure}[H]
\centering
\includegraphics[width=0.74\linewidth]{mem_modelo}
\caption{Modelo discretizado: malla, carga puntual $P=\valP$ aplicada hacia abajo en el nodo 7 y apoyos empotrados en los nodos 1, 3 y 6.}
\label{fig:mem-modelo}
\end{figure}

\noindent
% El reparto es 0,34 / 0,62 y no mitad y mitad: la tabla de nodos son tres
% columnas de numeros (3,5 cm de ancho natural) y la de la derecha carga ademas
% el bloque de parametros globales, cuya fila mas larga -"Apoyos / Nodos 1, 3, 6
% (empotrados)"- pide 9,4 cm y antes se salia al margen derecho.
\begin{minipage}[t]{0.34\linewidth}
\centering
\captionof{table}{Nodos.}
\label{tab:mem-nodos}
\begin{tabular}{c r r}
\toprule
Nodo & $x$ & $y$ \\
\midrule
\nodosRows \\
\bottomrule
\end{tabular}
\end{minipage}
\hfill
\begin{minipage}[t]{0.62\linewidth}
\centering
\captionof{table}{Elementos (conectividad y espesor).}
\label{tab:mem-elem}
\begin{tabular}{c c c c c c}
\toprule
Elem. & $N_1$ & $N_2$ & $N_3$ & $N_4$ & $t$ \\
\midrule
\elementosRows \\
\bottomrule
\end{tabular}

\vspace{0.9em}
\captionof{table}{Parámetros globales del análisis.}
\label{tab:mem-param}
\begin{tabular}{@{}l l@{}}
\toprule
\textbf{Parámetro} & \textbf{Valor} \\
\midrule
Módulo de elasticidad & $E = \valE$ \\
Coeficiente de Poisson & $\nu = \valNu$ \\
Espesor & $t = \valT$ \\
Análisis & Tensión plana \\
Carga & $F_y = -\valP$ (nodo 7) \\
Apoyos & Nodos 1, 3, 6 (empotrados) \\
\bottomrule
\end{tabular}
\end{minipage}

\section{Procedimiento}

\subsection{Matriz constitutiva}

La matriz constitutiva de tensión plana es la de la \autoref{eq:D-tension-plana}. Con $E=\valE$ y $\nu=\valNu$ resulta, idéntica para los cuatro elementos, que comparten material:
% Partida en dos renglones: las tres matrices 3x3 seguidas mas la fraccion no
% entran en el ancho de linea y se metian 0,9 cm en el margen derecho.
\begin{align*}
\bm{D}&=\frac{E}{1-\nu^{2}}
\begin{bmatrix} 1 & \nu & 0 \\ \nu & 1 & 0 \\ 0 & 0 & \tfrac{1-\nu}{2} \end{bmatrix}
=\valDfactor\;\DmatNorm \\[0.4em]
&=\Dmat .
\end{align*}

\subsection{Rigidez del elemento}

Se desarrolla el elemento E\showId{} (nodos \showNodes), el de mayor energía de deformación $U_e=\tfrac12\,\bm{u}_e^{\mathsf{T}}\bm{k}_e\,\bm{u}_e$ (\energiaInline); los tres restantes se calculan igual. Sus coordenadas nodales son
\[
\bm{X}_e=\Xemat .
\]
En el punto de Gauss $\gpA$, las derivadas de las funciones de forma de la \autoref{eq:N-q4} en coordenadas naturales valen
\[
\partial\bm{N}_{\xi\eta}=\dNnatA .
\]
El Jacobiano de la \autoref{eq:jacobiano}, su determinante (\autoref{eq:detJ}) y su inversa por cofactores (\autoref{eq:invJ}) son
% Idem: las tres igualdades en un renglon se pasaban del margen.
\begin{align*}
\bm{J}&=\partial\bm{N}_{\xi\eta}\,\bm{X}_e=\Jmata, \qquad
\det\bm{J}=\detJa, \\[0.4em]
\bm{J}^{-1}&=\invJa,
\end{align*}
y las derivadas físicas, obtenidas por la regla de la cadena de la \autoref{eq:fisica-natural}, y la matriz $\bm{B}$ de la \autoref{eq:matriz-B} ($3\times 8$, que ordena esas derivadas según $\bm{\varepsilon}=\bm{B}\,\bm{u}_e$) resultan
\[
\partial\bm{N}_{xy}=\bm{J}^{-1}\,\partial\bm{N}_{\xi\eta}=\dNphysA ,
\]
{\footnotesize\setlength{\arraycolsep}{4pt}
\[
\bm{B}=\Bmata .
\]
}%
La rigidez del elemento es la cuadratura de la \autoref{eq:ke-gauss} sobre los cuatro puntos de Gauss (pesos $w_g=1$; determinantes $\det\bm{J}_g \in \{\detJinline\}$):
\[
\bm{k}_e=\sum_{g=1}^{4} w_g\,\bm{B}_g^{\mathsf{T}}\,\bm{D}\,\bm{B}_g\,\lvert\det\bm{J}_g\rvert\,t .
\]
{\footnotesize\setlength{\arraycolsep}{3.5pt}
\[
\bm{k}_e=\keScale\cdot\keMat .
\]
}%

\subsection{Ensamblaje}

Cada $\bm{k}_e$ se dispersa a la matriz global según los grados de libertad (GDL) de sus nodos, con el operador de ensamblaje de la \autoref{eq:ensamblaje}; las cargas nodales arman $\bm{F}$:
\[
\bm{K}=\Ensamblaje_{e=1}^{4}\bm{k}_e \quad(\nDof\times\nDof),
\qquad
\bm{F}:\; F_{(7,y)}=-\valP,\ \text{resto } 0 .
\]
La \autoref{fig:mem-K} muestra la estructura de no-nulos de $\bm{K}$.

\begin{figure}[H]
\centering
\includegraphics[width=0.42\linewidth]{mem_Ksparsity}
\caption{Patrón de no-nulos de la matriz de rigidez global $\bm{K}$ ($\nDof\times\nDof$). Cada bloque $2\times 2$ acopla dos nodos vecinos; el nodo central 5, compartido por los cuatro elementos, concentra el mayor número de acoplamientos.}
\label{fig:mem-K}
\end{figure}

\subsection{Restricciones y solución}

Los nodos 1, 3 y 6 fijan sus dos GDL: quedan \nRestr{} GDL restringidos y \nFree{} libres. Reordenando el sistema en bloques libre ($f$) y restringido ($r$), según la partición de la \autoref{eq:particion},
\[
\begin{bmatrix} \bm{K}_{ff} & \bm{K}_{fr} \\ \bm{K}_{rf} & \bm{K}_{rr} \end{bmatrix}
\begin{bmatrix} \bm{u}_{f} \\ \bm{u}_{r} \end{bmatrix}
=\begin{bmatrix} \bm{F}_{f} \\ \bm{F}_{r} \end{bmatrix},
\]
con $\bm{u}_r=\bm{0}$ el subsistema libre ($\nFree\times\nFree$) se reduce a
\[
\bm{K}_{ff}\,\bm{u}_{f}=\bm{F}_{f},
\qquad
\bm{F}_{f}^{\mathsf{T}}=\bigl[\,0\ \cdots\ -\valP\ \cdots\ 0\,\bigr] .
\]
Su factorización entrega los desplazamientos nodales (\autoref{tab:mem-resultados}).

\subsection{Reacciones}

En los GDL restringidos, $\bm{R}=\bm{K}\,\bm{u}-\bm{F}$:
\begin{center}
\begin{tabular}{c r r}
\toprule
Nodo & $R_x$ & $R_y$ \\
\midrule
\reaccRows \\
\midrule
$\sum$ & $\RxTot$ & $\RyTot$ \\
\bottomrule
\end{tabular}
\end{center}
La suma vertical $\RyTot$ equilibra exactamente la carga aplicada $-\valP$, y la horizontal es nula: el modelo está en equilibrio.

\subsection{Recuperación de tensiones}

En cada punto de Gauss del elemento, $\bm{\varepsilon}=\bm{B}\,\bm{u}_e$ y $\bm{\sigma}=\bm{D}\,\bm{\varepsilon}$. En $\gpA$:
\[
\bm{\varepsilon}=\epsA,
\qquad
\bm{\sigma}=\sigA .
\]
Las tensiones principales y la de von Mises de la \autoref{eq:principales-vm}, en su forma plana ($\sigma_z = 0$ en tensión plana), resultan en ese punto
\[
\sigma_{1,2}=\frac{\sigma_x+\sigma_y}{2}\pm\sqrt{\Bigl(\tfrac{\sigma_x-\sigma_y}{2}\Bigr)^{2}+\tau_{xy}^{2}}=\savgA\pm\radiusA=\sigOneA,\ \sigTwoA ,
\]
\[
\sigma_{\mathrm{VM}}=\sqrt{\sigma_1^{2}-\sigma_1\sigma_2+\sigma_2^{2}}=\vmA .
\]
Repetida la operación en los cuatro puntos de Gauss —numerados $(-,-)$, $(-,+)$, $(+,-)$, $(+,+)$—, cada componente cartesiana se lleva a los nodos mediante la matriz de extrapolación constante $\bm{E}=\bm{M}^{-1}$ de la \autoref{eq:extrap-q4}, $\bm{\sigma}_{\mathrm{nod}}=\bm{E}\,\bm{\sigma}_{\mathrm{PG}}$; por ejemplo, para $\sigma_x$ (nodos \showNodes):
{\footnotesize\setlength{\arraycolsep}{4pt}
\[
\sigxNodal=\Eextrap\,\sigxGauss .
\]
}%
Las Tablas \ref{tab:mem-gauss} y \ref{tab:mem-extrap} reúnen las cuatro componentes en los puntos de Gauss y, ya extrapoladas, en los nodos del elemento E\showId{}. La extrapolación se aplica únicamente a las componentes cartesianas: la tensión de von Mises de cada nodo se recalcula con la \autoref{eq:principales-vm} a partir de las componentes ya extrapoladas, por ser una función no lineal de ellas.

\noindent
\begin{minipage}[t]{0.49\linewidth}
\centering
\captionof{table}{Tensiones en los puntos de Gauss (E\showId).}
\label{tab:mem-gauss}
\footnotesize
\begin{tabular}{c r r r r}
\toprule
PG & $\sigma_x$ & $\sigma_y$ & $\tau_{xy}$ & $\sigma_{\mathrm{VM}}$ \\
\midrule
\gaussStressRows \\
\bottomrule
\end{tabular}
\end{minipage}
\hfill
\begin{minipage}[t]{0.49\linewidth}
\centering
\captionof{table}{Tensiones extrapoladas a los nodos (E\showId).}
\label{tab:mem-extrap}
\footnotesize
\begin{tabular}{c r r r r}
\toprule
Nodo & $\sigma_x$ & $\sigma_y$ & $\tau_{xy}$ & $\sigma_{\mathrm{VM}}$ \\
\midrule
\nodalStressShowRows \\
\bottomrule
\end{tabular}
\end{minipage}

\medskip
\noindent El campo nodal final (\autoref{tab:mem-resultados}) promedia, en cada nodo, las contribuciones de los elementos que lo comparten. El promedio recae sobre las componentes cartesianas; las tensiones principales y la de von Mises del nodo se obtienen después, a partir de esas componentes promediadas, de modo que cada fila de la tabla es consistente consigo misma.

\section{Resultados}

La \autoref{tab:mem-resultados} reúne, por nodo, los desplazamientos (solución del sistema) y las tensiones promediadas (recuperadas y promediadas entre elementos).

\begin{table}[H]
\centering
\caption{Resultados nodales: desplazamientos y tensiones promediadas.}
\label{tab:mem-resultados}
\setlength{\tabcolsep}{5pt}
\begin{tabular}{c r r r r r r}
\toprule
Nodo & $u_x$ & $u_y$ & $\sigma_x$ & $\sigma_y$ & $\tau_{xy}$ & $\sigma_{\mathrm{VM}}$ \\
\midrule
\nodalAllRows \\
\bottomrule
\end{tabular}
\end{table}

\noindent El desplazamiento máximo es $\lvert\bm{u}\rvert_{\max}=\umaxMag$ en el nodo \umaxNode{} (donde se aplica la carga), y la tensión de von Mises máxima, $\sigma_{\mathrm{VM},\max}=\vmMax$ en el mismo nodo. Las figuras \ref{fig:mem-deformada} y \ref{fig:mem-contorno} muestran la configuración deformada y el contorno de von Mises sobre la malla deformada, tal como los presenta el post-proceso del software.

% Dos figuras lado a lado SIN entorno figure: con \floatstyle{plaintop} un
% float con dos \caption conservaba solo el ultimo rotulo (la G.3 quedaba sin
% rotulo y el de la G.4 caia sobre la imagen equivocada). \captionof numera y
% ancla cada una por separado; el rotulo va arriba, como pide APA.
\medskip
\noindent
\begin{minipage}[t]{0.49\linewidth}
  \centering
  \captionof{figure}{Configuración deformada (amplificada) sobre la malla original.}
  \label{fig:mem-deformada}
  \includegraphics[width=\linewidth]{mem_deformada}
\end{minipage}
\hfill
\begin{minipage}[t]{0.49\linewidth}
  \centering
  \captionof{figure}{Contorno de tensión de von Mises sobre la malla deformada.}
  \label{fig:mem-contorno}
  \includegraphics[width=\linewidth]{mem_contorno_vm}
\end{minipage}
